\subsubsection{Subproceso 3.A: Verificaci\'on y Control de Reglas Acad\'emicas}

El subproceso se activa cuando el estudiante concluye el armado de su horario en el Subproceso~2 y el sistema genera la Pre-Ficha de Matr\'icula Registrada, la cual bloquea temporalmente los cupos elegidos en la base de datos del SIGA-ORCE. En este tramo del proceso, la propuesta acad\'emica entra en una etapa de auditor\'ia e inspecci\'on para asegurar que se ajuste estrictamente a las normas de la universidad antes de su asentamiento definitivo.

En primer lugar, el SIGA-ORCE efect\'ua una verificaci\'on algor\'itmica autom\'atica: valida la matriz de prerrequisitos, corrobora que no se supere el tope de cr\'editos correspondiente al promedio ponderado del alumno e impide cualquier cruce de horario en asignaturas de car\'acter pr\'actico o de laboratorio. Asimismo, comprueba que los estudiantes en condici\'on de riesgo acad\'emico (Grupo Cero) hayan inscrito obligatoriamente sus asignaturas desaprobadas del ciclo m\'as atrasado, conforme lo estipulan el Reglamento de Matr\'icula (Arts.~8\textdegree{} y 48\textdegree{})~\cite{reglamento2022} y el Estatuto de la UNI (Art.~253\textdegree{})~\cite{ley30220}. Aquellas fichas que satisfacen estas reglas obtienen la condici\'on de aprobadas preliminarmente.

En segundo lugar, se abre el periodo de rectificaciones para atender las observaciones surgidas o las contingencias del estudiante (como secciones agotadas en su turno o cruces de horarios te\'oricos). El alumno formula su requerimiento mediante el formulario digital institucional indicando si solicita inclusi\'on, retiro o cambio de secci\'on. La Direcci\'on de Escuela Profesional, apoyada por la Comisi\'on de Matr\'icula de la Facultad, analiza cada pedido. Si se autoriza un cruce te\'orico permitido (hasta un m\'aximo de cuatro horas pedag\'ogicas) o se concede una vacante, la Escuela aprueba la solicitud e ingresa manualmente las modificaciones al SIGA, generando la Pre-Ficha Subsanada. Por el contrario, si el estudiante no subsana una observaci\'on cr\'itica dentro de la semana de depuraci\'on, el sistema cancela la reserva temporal y libera las vacantes. El subproceso termina cuando la Pre-Ficha Aprobada o Subsanada queda emitida y lista para la confirmaci\'on oficial en el Subproceso~3.B.

Durante el levantamiento de informaci\'on del periodo 2026-2 se registraron las siguientes dificultades operativas:
\begin{enumerate}
    \item \textbf{Cuello de botella en la atenci\'on de solicitudes:} La Direcci\'on de Escuela experimenta una sobrecarga masiva de pedidos de rectificaci\'on que tardan entre 48 y 72 horas en procesarse. Gran parte de estos reclamos se deben a cruces de horario y vacantes agotadas que la plataforma no previene de forma autom\'atica en el paso anterior.
    \item \textbf{Revisi\'on manual de reincorporados y convalidados:} Para estudiantes que regresan de suspensi\'on o tienen convalidaciones recientes, el sistema SIGA no actualiza sus reglas autom\'aticamente, lo que obliga a la Escuela a cotejar sus expedientes en papel o archivos aislados caso por caso.
    \item \textbf{Efecto de arrastre hacia el cierre oficial:} La demora en resolver las Pre-Fichas observadas traslada retrasos cr\'iticos al Subproceso~3.B. Como consecuencia, el cierre formal de matr\'icula se posterga y diversos estudiantes acuden a clases durante la primera semana sin figurar en los padrones provisionales de los docentes.
\end{enumerate}

Esta descripci\'on del proceso actual complementa el diagn\'ostico de causas y sustenta las mejoras de automatizaci\'on y control que se formulan en las siguientes secciones.
